\documentclass[11pt]{article}
\usepackage[margin=0.8in]{geometry}
\usepackage{amsmath,microtype}
\newcommand{\guideheading}[1]{\par\medskip\noindent{\large\bfseries #1}\par\smallskip}
\begin{document}
\begin{center}
{\Large Guide: Why Local Cycle Certificates Do Not Cover Everything}\par\medskip
Collatz proof project\quad---\quad September 5, 2026
\end{center}
\guideheading{The question}
The inverse excursion can repair either cycle mismatch once both base
partners are known to converge. Could we now prove every parameter by
covering it with a certificate built from a strictly smaller base?
Lean identifies an infinite family that this particular bootstrap misses.

\guideheading{The basic speed limit}
A shortcut Collatz step cannot shrink a number faster than dividing it by
two. Therefore, if a target reaches 1 or 2 within $d$ steps, its starting
value is at most $2^{d+1}$.

A base parameter $r$ has target $27r+20$. To use a certificate whose binary
precision is $d$, with the base target already in the cycle by that depth,
we must have
\[
 27r+20\leq2^{d+1}.
\]
This forces the base to be small compared with the binary modulus.

\guideheading{The missed family}
Consider parameters $u=2^k-1$, whose binary digits are all ones.
Every shorter binary prefix has the same form: $r=2^d-1$.
That base is far too large to satisfy the required speed limit.

For example, parameter 31 has proper prefix values 1, 3, 7, and 15.
None has its target in the cycle within its corresponding prefix depth.
At a larger precision the base becomes 31 itself, so it is no longer
strictly smaller. Lean proves the analogous statement for every exponent
$k$ and every choice of depth $d$.

\guideheading{Why more inverse excursions do not fix this}
The obstruction appears before choosing excursions: the required smaller
base does not meet the cycle-entry condition. Adding more bridges after
the cycle, or allowing the inverse bridge, cannot change that fact.

This explains why a successful constructor for many individually verified
base pairs is not yet a proof that every new parameter is covered.

\guideheading{What the theorem does not say}
It does not say these parameters fail to converge. It does not exclude a
direct merge before reaching the cycle, a different induction, or another
family of certificates. The inverse bridge theorem remains valid with
its explicit premises.

The result rules out one attempted use of those local certificates as a
universal bootstrap. A complete proof needs an additional argument that
does not depend on selecting a smaller base satisfying this condition.
Collatz remains unproved.

\guideheading{Verification}
Build \texttt{Collatz.CycleCylinderBoundary} from the repository's
\texttt{lean/} directory. Its five public declarations are included in the
axiom audit. The technical paper states the exact hypotheses and their
all-depth contradiction.
\end{document}
